% GENERATED FILE. Edit canonical lessons, then run npm run generate:books.
% canonical-source-hash: fnv1a64:644f06b3ccbc457a
% canonical-lessons: MR-R14-mitra-r2, MR-R14-kutumb-r2, MR-R14-siblings, MR-R14-eye-ear, MR-R14-mouth-nose, MR-R14-food, MR-R14-doing, MR-R14-mind, MR-R14-core, MR-R14-mitra-r3, MR-R14-numbers, MR-R14-kutumb-r3

\chapter{Family Memory at Real Distance}
\label{ch:mr-family-memory}
\hlchaptermodality{21}

\begin{chapteropening}
\textbf{By the end of this chapter:} \emph{I can retrieve \mr{मित्र} and \mr{कुटुंब} after short, medium, and durable gaps; explain the binding root and tatsama path of \mr{मित्र}; and recall that \mr{कुटुंब} is neuter in Marathi.}

After genuine intervening retrievals, produce the five-part comparison: \mr{मित्र}, its binding root and tatsama path, then \mr{कुटुंब} and its Marathi neuter gender.
\end{chapteropening}

\section[\texorpdfstring{\mr{मित्र} (mitra)}{मित्र (mitra)}]{After a gap — \mr{मित्र}}

\label{lesson:MR-R14-mitra-r2}



\begin{quote}
Friend. Old root meaning \textbf{to bind}. Borrowed whole from Sanskrit. Say the three answers: \textbf{\mr{मित्र} — bind — tatsama}.
\end{quote}

\subsection*{Your turn}
Answer once, then repair only the missing piece.

\begin{tcolorbox}[breakable,colback=teal!4,colframe=teal!35!black,title={Before you move on}]
Say \textbf{\mr{मित्र}} once more, then stop.
\end{tcolorbox}
\section[\texorpdfstring{\mr{कुटुंब} (kutumb)}{कुटुंब (kutumb)}]{After a gap — \mr{कुटुंब}}

\label{lesson:MR-R14-kutumb-r2}



\begin{quote}
Family. Marathi gender. Answer: \textbf{\mr{कुटुंब} — neuter}.
\end{quote}

\subsection*{Your turn}
Answer once, then repair only the missing piece.

\begin{tcolorbox}[breakable,colback=teal!4,colframe=teal!35!black,title={Before you move on}]
Say \textbf{\mr{कुटुंब}} once more.
\end{tcolorbox}
\section[\texorpdfstring{\mr{भाऊ} / \mr{बहीण} (bhau / bahin)}{भाऊ / बहीण (bhau / bahin)}]{Sibling retrieval}

\label{lesson:MR-R14-siblings}



\begin{quote}
Say \textbf{brother, sister}: \textbf{\mr{भाऊ}, \mr{बहीण}}.
\end{quote}

\subsection*{The two words and their histories}
\textbf{\mr{भाऊ}} is a secure cousin of English \textbf{brother}. Do not extend that claim to \textbf{\mr{बहीण}} and English \emph{sister}; those roots are unrelated.

\subsection*{Your turn}
Answer once, then repair only the missing piece.

\begin{tcolorbox}[breakable,colback=teal!4,colframe=teal!35!black,title={Before you move on}]
\textbf{\mr{भाऊ}, \mr{बहीण}.}
\end{tcolorbox}
\section[\texorpdfstring{\mr{डोळा} / \mr{कान} (dola / kaan)}{डोळा / कान (dola / kaan)}]{Eye and ear retrieval}

\label{lesson:MR-R14-eye-ear}



\begin{quote}
Eye: \textbf{\mr{डोळा}}. Ear: \textbf{\mr{कान}}. In \textbf{\mr{डोळा}}, keep \textbf{\mr{ळ}} retroflex; do not replace it with ordinary \textbf{\mr{ल}}.
\end{quote}

\subsection*{Your turn}
Answer once, then repair only the missing piece.

\begin{tcolorbox}[breakable,colback=teal!4,colframe=teal!35!black,title={Before you move on}]
\textbf{\mr{डोळा}, \mr{कान}.}
\end{tcolorbox}
\section[\texorpdfstring{\mr{तोंड} / \mr{नाक} (tond / naak)}{तोंड / नाक (tond / naak)}]{Mouth and nose retrieval}

\label{lesson:MR-R14-mouth-nose}



\begin{quote}
Everyday \textbf{mouth} is \textbf{\mr{तोंड}}; formal \textbf{\mr{मुख}} stands beside it. \textbf{Nose} is \textbf{\mr{नाक}}.
\end{quote}

\subsection*{Your turn}
Answer once, then repair only the missing piece.

\begin{tcolorbox}[breakable,colback=teal!4,colframe=teal!35!black,title={Before you move on}]
\textbf{\mr{तोंड}, \mr{नाक}.}
\end{tcolorbox}
\section[\texorpdfstring{\mr{पाणी} / \mr{चहा} / \mr{दूध} (paani / chaha / dudh)}{पाणी / चहा / दूध (paani / chaha / dudh)}]{Three request nouns}

\label{lesson:MR-R14-food}



\begin{quote}
Say \textbf{water, tea, milk} in order: \textbf{\mr{पाणी}, \mr{चहा}, \mr{दूध}}.
\end{quote}

\subsection*{Your turn}
Answer once, then repair only the missing piece.

\begin{tcolorbox}[breakable,colback=teal!4,colframe=teal!35!black,title={Before you move on}]
Repeat once, then stop.
\end{tcolorbox}
\section[\texorpdfstring{\mr{घेणे} / \mr{विचारणे} … (ghene / vicharne …)}{घेणे / विचारणे … (ghene / vicharne …)}]{Three doing verbs}

\label{lesson:MR-R14-doing}



\begin{quote}
Take: \textbf{\mr{घेणे}}. Ask: \textbf{\mr{विचारणे}}. Help: \textbf{\mr{मदत} \mr{करणे}}.
\end{quote}

\subsection*{Your turn}
Answer once, then repair only the missing piece.

\begin{tcolorbox}[breakable,colback=teal!4,colframe=teal!35!black,title={Before you move on}]
Repeat the trio once.
\end{tcolorbox}
\section[\texorpdfstring{\mr{विचार} \mr{करणे} / \mr{समजणे} … (vichar karne …)}{विचार करणे / समजणे … (vichar karne …)}]{Three mind-and-page verbs}

\label{lesson:MR-R14-mind}



\begin{quote}
Think: \textbf{\mr{विचार} \mr{करणे}}. Understand: \textbf{\mr{समजणे}}. Write: \textbf{\mr{लिहिणे}}.
\end{quote}

\subsection*{Your turn}
Answer once, then repair only the missing piece.

\begin{tcolorbox}[breakable,colback=teal!4,colframe=teal!35!black,title={Before you move on}]
Repeat once.
\end{tcolorbox}
\section[\texorpdfstring{\mr{असणे} / \mr{जाणे} / \mr{येणे} (asne / jane / yene)}{असणे / जाणे / येणे (asne / jane / yene)}]{Three core verbs}

\label{lesson:MR-R14-core}



\begin{quote}
Be: \textbf{\mr{असणे}}. Go: \textbf{\mr{जाणे}}. Come: \textbf{\mr{येणे}}.
\end{quote}

\subsection*{Your turn}
Answer once, then repair only the missing piece.

\begin{tcolorbox}[breakable,colback=teal!4,colframe=teal!35!black,title={Before you move on}]
Repeat once.
\end{tcolorbox}
\section[\texorpdfstring{\mr{मित्र} (mitra)}{मित्र (mitra)}]{Durable retrieval — \mr{मित्र}}

\label{lesson:MR-R14-mitra-r3}



\begin{quote}
Recover the word for \textbf{friend}, its old root meaning, and the name for a Sanskrit form borrowed whole.
\end{quote}

\subsection*{Your turn}
Answer once, then repair only the missing piece.

\begin{tcolorbox}[breakable,colback=teal!4,colframe=teal!35!black,title={Before you move on}]
\textbf{\mr{मित्र}}. One clean retrieval is enough.
\end{tcolorbox}
\section[\texorpdfstring{\mr{एक} \mr{ते} \mr{पाच} (ek te paach)}{एक ते पाच (ek te paach)}]{Count, then notice}

\label{lesson:MR-R14-numbers}



\begin{quote}
Count one through five. Then say \textbf{\mr{चार}} with Marathi \textbf{tsā-}, and remember that a language can innovate in one place while retaining an older feature in another.
\end{quote}

\subsection*{Your turn}
Answer once, then repair only the missing piece.

\begin{tcolorbox}[breakable,colback=teal!4,colframe=teal!35!black,title={Before you move on}]
Count once more without looking.
\end{tcolorbox}
\section[\texorpdfstring{\mr{मित्र} / \mr{कुटुंब} (mitra / kutumb)}{मित्र / कुटुंब (mitra / kutumb)}]{Durable family comparison}

\label{lesson:MR-R14-kutumb-r3}



\begin{quote}
Friend: \textbf{\mr{मित्र}}; old root \textbf{to bind}; \textbf{tatsama}. Family: \textbf{\mr{कुटुंब}}; \textbf{neuter} in Marathi.
\end{quote}

\subsection*{Your turn}
Answer once, then repair only the missing piece.

\begin{tcolorbox}[breakable,colback=teal!4,colframe=teal!35!black,title={Before you move on}]
\textbf{\mr{मित्र}, \mr{कुटुंब}.} End the session here.
\end{tcolorbox}
